\section{Founding Principles}


The architecture of Lux is guided by eight principles. These are not
marketing statements---they are constraints on the solution space,
each of which can veto any proposed change to the network.

\begin{enumerate}[leftmargin=*,nosep]
  \item \textbf{Approachability}: Anyone should be able to build on Lux.
        APIs and libraries target multiple languages (Go, TypeScript,
        Python) with identical semantics.
  \item \textbf{Availability}: A financial system that is intermittently
        available is not a financial system. Lux targets five-nines
        RPC availability through multi-region validators and
        leaderless consensus.
  \item \textbf{Composability}: Small, modular primitives that
        assemble into larger systems. Every chain, VM, and precompile
        is usable independently or in combination.
  \item \textbf{Flexibility}: Well-defined extension points (custom
        VMs, pluggable consensus, per-chain genesis) with exactly one
        way to use each.
  \item \textbf{Longevity}: Bets that will still be correct in a
        decade. Lattice cryptography over RSA. UTXO where quantum
        resistance matters. Leaderless over leader-based consensus.
  \item \textbf{Interoperability}: A chain that cannot communicate
        with other chains is a database with extra steps. Native
        cross-chain messaging and threshold-signed bridges.
  \item \textbf{Security}: Defense in depth. Post-quantum signatures
        at consensus. Threshold cryptography for key management.
        Formal verification of the core protocol.
  \item \textbf{Scalability}: Horizontal chain creation, parallelized
        EVM execution, and a DAG-based consensus structure that does
        not serialize all transactions through a single leader.
\end{enumerate}

These principles are developed in depth in a companion document, the
\emph{Ethos of Lux}~\cite{lux-ethos}.

%% ============================================================
